\newcommand{\coursemodule}{Module 09}
% ════════════════════════════════════════════════════════════════════════════
% Pixels, Frames & Transformers: shared Beamer preamble (WUR style)
% Adapted from the preamble of Bags, Vectors & Transformers.
%
% Use in a module deck (from <part>/<module>/slides/):
%   \newcommand{\coursemodule}{Module 01}
%   % ════════════════════════════════════════════════════════════════════════════
% Pixels, Frames & Transformers: shared Beamer preamble (WUR style)
% Adapted from the preamble of Bags, Vectors & Transformers.
%
% Use in a module deck (from <part>/<module>/slides/):
%   \newcommand{\coursemodule}{Module 01}
%   % ════════════════════════════════════════════════════════════════════════════
% Pixels, Frames & Transformers: shared Beamer preamble (WUR style)
% Adapted from the preamble of Bags, Vectors & Transformers.
%
% Use in a module deck (from <part>/<module>/slides/):
%   \newcommand{\coursemodule}{Module 01}
%   \input{../../../shared/preamble.tex}
%   \begin{document}
%   \moduletitle{Images as Data}
%   ...
% Compile with: latexmk -pdf <deck>.tex
% ════════════════════════════════════════════════════════════════════════════

\documentclass[aspectratio=169]{beamer}

% ── Packages ────────────────────────────────────────────────────────────────
\usepackage[T1]{fontenc}
\usepackage[utf8]{inputenc}
\usepackage{lmodern}
\usepackage{microtype}
\usepackage{xcolor}
\usepackage{tikz}
\usetikzlibrary{arrows.meta, positioning, calc, shapes.geometric, decorations.pathreplacing, matrix, fit}
\usepackage{booktabs}
\usepackage{fontawesome5}
\usepackage{array}
\usepackage[most]{tcolorbox}

\providecommand{\coursemodule}{}

% ── WUR color palette ───────────────────────────────────────────────────────
\definecolor{wurgreen}{HTML}{34B233}
\definecolor{darkgreen}{HTML}{1F6B1E}
\definecolor{lightgreen}{HTML}{E8F7E8}
\definecolor{darknavy}{HTML}{1A1A2E}
\definecolor{midgray}{HTML}{6B7280}
\definecolor{lightgray}{HTML}{F3F4F6}
\definecolor{bodytext}{HTML}{1F2937}
\definecolor{codegreen}{HTML}{9AE6B4}
\definecolor{warnred}{HTML}{FEE2E2}
\definecolor{warnborder}{HTML}{EF4444}

% ── Beamer theme ────────────────────────────────────────────────────────────
\usetheme{default}
\usecolortheme{default}
\usefonttheme{professionalfonts}
\setbeamertemplate{navigation symbols}{}
\setbeamertemplate{itemize items}{\textcolor{wurgreen}{\textbullet}}
\setbeamertemplate{itemize subitem}{\textcolor{midgray}{--}}
\setbeamertemplate{enumerate items}{\textcolor{wurgreen}{\textbf{\insertenumlabel.}}}
\setbeamercolor{frametitle}{fg=white, bg=darknavy}
\setbeamertemplate{frametitle}{%
  \begin{beamercolorbox}[wd=\paperwidth, ht=1.1cm, dp=0.2cm, leftskip=0.6cm]{frametitle}
    \usebeamerfont{frametitle}\insertframetitle
  \end{beamercolorbox}%
}
\setbeamerfont{frametitle}{size=\large, series=\bfseries}
\setbeamercolor{background canvas}{bg=white}
\setbeamercolor{normal text}{fg=bodytext}
\setbeamerfont{normal text}{size=\small}
\setbeamertemplate{footline}{%
  \begin{beamercolorbox}[wd=\paperwidth, ht=0.4cm, dp=0.15cm]{footline}%
    \hspace{0.5cm}%
    \textcolor{midgray}{\tiny Pixels, Frames \& Transformers · \coursemodule\ · Wageningen University \& Research}%
    \hfill%
    \textcolor{midgray}{\tiny \insertframenumber}%
    \hspace{0.4cm}%
  \end{beamercolorbox}%
}

% ── Boxes ───────────────────────────────────────────────────────────────────
\tcbset{wurbase/.style={arc=3pt, boxrule=0.6pt, left=7pt, right=7pt, top=5pt, bottom=5pt,
                        fontupper=\small, before skip=4pt, after skip=4pt}}
\newtcolorbox{greenbox}[1][]{wurbase, colback=lightgreen, colframe=wurgreen, #1}
\newtcolorbox{graybox}[1][]{wurbase, colback=lightgray, colframe=lightgray, #1}
\newtcolorbox{navybox}[1][]{wurbase, colback=darknavy, colframe=darknavy, colupper=white, #1}
\newtcolorbox{warnbox}[1][]{wurbase, colback=warnred, colframe=warnborder, #1}
\newtcolorbox{codebox}[1][]{wurbase, colback=darknavy, colframe=darknavy, colupper=white,
                            fontupper=\ttfamily\scriptsize, left=6pt, top=4pt, bottom=4pt, #1}

% ── Title slide ─────────────────────────────────────────────────────────────
\newcommand{\moduletitle}[1]{%
{
\setbeamertemplate{footline}{}
\begin{frame}[plain]
  \begin{tikzpicture}[remember picture, overlay]
    \fill[darknavy] (current page.south west) rectangle (current page.north east);
    \fill[wurgreen] (current page.south west) rectangle ([xshift=0.45cm]current page.north west);
    \node[anchor=west, text=white, font=\huge\bfseries, text width=13cm]
      at ([xshift=1.1cm, yshift=1.2cm]current page.west) {Pixels, Frames \& Transformers};
    \node[anchor=west, text=wurgreen, font=\large, text width=13cm]
      at ([xshift=1.1cm, yshift=0.05cm]current page.west) {#1 · \coursemodule};
    \node[anchor=west, text=midgray, font=\small, text width=12cm]
      at ([xshift=1.1cm, yshift=-1.5cm]current page.west)
      {Denise J. Roth, PhD Candidate\\
       Strategic Communication Group · Wageningen University \& Research};
  \end{tikzpicture}
\end{frame}
}}

% ── Small helpers ───────────────────────────────────────────────────────────
\newcommand{\keyterm}[1]{\textbf{\textcolor{darkgreen}{#1}}}
\newcommand{\muted}[1]{\textcolor{midgray}{#1}}
\newcommand{\credit}[1]{{\tiny\textcolor{midgray}{#1}}}
\newcommand{\notebookslide}[1]{%
\begin{frame}{In the Notebook}
\vspace{0.1cm}
#1
\end{frame}}
\newcommand{\nbitem}[2]{\item[\textcolor{wurgreen}{#1}] #2}

%   \begin{document}
%   \moduletitle{Images as Data}
%   ...
% Compile with: latexmk -pdf <deck>.tex
% ════════════════════════════════════════════════════════════════════════════

\documentclass[aspectratio=169]{beamer}

% ── Packages ────────────────────────────────────────────────────────────────
\usepackage[T1]{fontenc}
\usepackage[utf8]{inputenc}
\usepackage{lmodern}
\usepackage{microtype}
\usepackage{xcolor}
\usepackage{tikz}
\usetikzlibrary{arrows.meta, positioning, calc, shapes.geometric, decorations.pathreplacing, matrix, fit}
\usepackage{booktabs}
\usepackage{fontawesome5}
\usepackage{array}
\usepackage[most]{tcolorbox}

\providecommand{\coursemodule}{}

% ── WUR color palette ───────────────────────────────────────────────────────
\definecolor{wurgreen}{HTML}{34B233}
\definecolor{darkgreen}{HTML}{1F6B1E}
\definecolor{lightgreen}{HTML}{E8F7E8}
\definecolor{darknavy}{HTML}{1A1A2E}
\definecolor{midgray}{HTML}{6B7280}
\definecolor{lightgray}{HTML}{F3F4F6}
\definecolor{bodytext}{HTML}{1F2937}
\definecolor{codegreen}{HTML}{9AE6B4}
\definecolor{warnred}{HTML}{FEE2E2}
\definecolor{warnborder}{HTML}{EF4444}

% ── Beamer theme ────────────────────────────────────────────────────────────
\usetheme{default}
\usecolortheme{default}
\usefonttheme{professionalfonts}
\setbeamertemplate{navigation symbols}{}
\setbeamertemplate{itemize items}{\textcolor{wurgreen}{\textbullet}}
\setbeamertemplate{itemize subitem}{\textcolor{midgray}{--}}
\setbeamertemplate{enumerate items}{\textcolor{wurgreen}{\textbf{\insertenumlabel.}}}
\setbeamercolor{frametitle}{fg=white, bg=darknavy}
\setbeamertemplate{frametitle}{%
  \begin{beamercolorbox}[wd=\paperwidth, ht=1.1cm, dp=0.2cm, leftskip=0.6cm]{frametitle}
    \usebeamerfont{frametitle}\insertframetitle
  \end{beamercolorbox}%
}
\setbeamerfont{frametitle}{size=\large, series=\bfseries}
\setbeamercolor{background canvas}{bg=white}
\setbeamercolor{normal text}{fg=bodytext}
\setbeamerfont{normal text}{size=\small}
\setbeamertemplate{footline}{%
  \begin{beamercolorbox}[wd=\paperwidth, ht=0.4cm, dp=0.15cm]{footline}%
    \hspace{0.5cm}%
    \textcolor{midgray}{\tiny Pixels, Frames \& Transformers · \coursemodule\ · Wageningen University \& Research}%
    \hfill%
    \textcolor{midgray}{\tiny \insertframenumber}%
    \hspace{0.4cm}%
  \end{beamercolorbox}%
}

% ── Boxes ───────────────────────────────────────────────────────────────────
\tcbset{wurbase/.style={arc=3pt, boxrule=0.6pt, left=7pt, right=7pt, top=5pt, bottom=5pt,
                        fontupper=\small, before skip=4pt, after skip=4pt}}
\newtcolorbox{greenbox}[1][]{wurbase, colback=lightgreen, colframe=wurgreen, #1}
\newtcolorbox{graybox}[1][]{wurbase, colback=lightgray, colframe=lightgray, #1}
\newtcolorbox{navybox}[1][]{wurbase, colback=darknavy, colframe=darknavy, colupper=white, #1}
\newtcolorbox{warnbox}[1][]{wurbase, colback=warnred, colframe=warnborder, #1}
\newtcolorbox{codebox}[1][]{wurbase, colback=darknavy, colframe=darknavy, colupper=white,
                            fontupper=\ttfamily\scriptsize, left=6pt, top=4pt, bottom=4pt, #1}

% ── Title slide ─────────────────────────────────────────────────────────────
\newcommand{\moduletitle}[1]{%
{
\setbeamertemplate{footline}{}
\begin{frame}[plain]
  \begin{tikzpicture}[remember picture, overlay]
    \fill[darknavy] (current page.south west) rectangle (current page.north east);
    \fill[wurgreen] (current page.south west) rectangle ([xshift=0.45cm]current page.north west);
    \node[anchor=west, text=white, font=\huge\bfseries, text width=13cm]
      at ([xshift=1.1cm, yshift=1.2cm]current page.west) {Pixels, Frames \& Transformers};
    \node[anchor=west, text=wurgreen, font=\large, text width=13cm]
      at ([xshift=1.1cm, yshift=0.05cm]current page.west) {#1 · \coursemodule};
    \node[anchor=west, text=midgray, font=\small, text width=12cm]
      at ([xshift=1.1cm, yshift=-1.5cm]current page.west)
      {Denise J. Roth, PhD Candidate\\
       Strategic Communication Group · Wageningen University \& Research};
  \end{tikzpicture}
\end{frame}
}}

% ── Small helpers ───────────────────────────────────────────────────────────
\newcommand{\keyterm}[1]{\textbf{\textcolor{darkgreen}{#1}}}
\newcommand{\muted}[1]{\textcolor{midgray}{#1}}
\newcommand{\credit}[1]{{\tiny\textcolor{midgray}{#1}}}
\newcommand{\notebookslide}[1]{%
\begin{frame}{In the Notebook}
\vspace{0.1cm}
#1
\end{frame}}
\newcommand{\nbitem}[2]{\item[\textcolor{wurgreen}{#1}] #2}

%   \begin{document}
%   \moduletitle{Images as Data}
%   ...
% Compile with: latexmk -pdf <deck>.tex
% ════════════════════════════════════════════════════════════════════════════

\documentclass[aspectratio=169]{beamer}

% ── Packages ────────────────────────────────────────────────────────────────
\usepackage[T1]{fontenc}
\usepackage[utf8]{inputenc}
\usepackage{lmodern}
\usepackage{microtype}
\usepackage{xcolor}
\usepackage{tikz}
\usetikzlibrary{arrows.meta, positioning, calc, shapes.geometric, decorations.pathreplacing, matrix, fit}
\usepackage{booktabs}
\usepackage{fontawesome5}
\usepackage{array}
\usepackage[most]{tcolorbox}

\providecommand{\coursemodule}{}

% ── WUR color palette ───────────────────────────────────────────────────────
\definecolor{wurgreen}{HTML}{34B233}
\definecolor{darkgreen}{HTML}{1F6B1E}
\definecolor{lightgreen}{HTML}{E8F7E8}
\definecolor{darknavy}{HTML}{1A1A2E}
\definecolor{midgray}{HTML}{6B7280}
\definecolor{lightgray}{HTML}{F3F4F6}
\definecolor{bodytext}{HTML}{1F2937}
\definecolor{codegreen}{HTML}{9AE6B4}
\definecolor{warnred}{HTML}{FEE2E2}
\definecolor{warnborder}{HTML}{EF4444}

% ── Beamer theme ────────────────────────────────────────────────────────────
\usetheme{default}
\usecolortheme{default}
\usefonttheme{professionalfonts}
\setbeamertemplate{navigation symbols}{}
\setbeamertemplate{itemize items}{\textcolor{wurgreen}{\textbullet}}
\setbeamertemplate{itemize subitem}{\textcolor{midgray}{--}}
\setbeamertemplate{enumerate items}{\textcolor{wurgreen}{\textbf{\insertenumlabel.}}}
\setbeamercolor{frametitle}{fg=white, bg=darknavy}
\setbeamertemplate{frametitle}{%
  \begin{beamercolorbox}[wd=\paperwidth, ht=1.1cm, dp=0.2cm, leftskip=0.6cm]{frametitle}
    \usebeamerfont{frametitle}\insertframetitle
  \end{beamercolorbox}%
}
\setbeamerfont{frametitle}{size=\large, series=\bfseries}
\setbeamercolor{background canvas}{bg=white}
\setbeamercolor{normal text}{fg=bodytext}
\setbeamerfont{normal text}{size=\small}
\setbeamertemplate{footline}{%
  \begin{beamercolorbox}[wd=\paperwidth, ht=0.4cm, dp=0.15cm]{footline}%
    \hspace{0.5cm}%
    \textcolor{midgray}{\tiny Pixels, Frames \& Transformers · \coursemodule\ · Wageningen University \& Research}%
    \hfill%
    \textcolor{midgray}{\tiny \insertframenumber}%
    \hspace{0.4cm}%
  \end{beamercolorbox}%
}

% ── Boxes ───────────────────────────────────────────────────────────────────
\tcbset{wurbase/.style={arc=3pt, boxrule=0.6pt, left=7pt, right=7pt, top=5pt, bottom=5pt,
                        fontupper=\small, before skip=4pt, after skip=4pt}}
\newtcolorbox{greenbox}[1][]{wurbase, colback=lightgreen, colframe=wurgreen, #1}
\newtcolorbox{graybox}[1][]{wurbase, colback=lightgray, colframe=lightgray, #1}
\newtcolorbox{navybox}[1][]{wurbase, colback=darknavy, colframe=darknavy, colupper=white, #1}
\newtcolorbox{warnbox}[1][]{wurbase, colback=warnred, colframe=warnborder, #1}
\newtcolorbox{codebox}[1][]{wurbase, colback=darknavy, colframe=darknavy, colupper=white,
                            fontupper=\ttfamily\scriptsize, left=6pt, top=4pt, bottom=4pt, #1}

% ── Title slide ─────────────────────────────────────────────────────────────
\newcommand{\moduletitle}[1]{%
{
\setbeamertemplate{footline}{}
\begin{frame}[plain]
  \begin{tikzpicture}[remember picture, overlay]
    \fill[darknavy] (current page.south west) rectangle (current page.north east);
    \fill[wurgreen] (current page.south west) rectangle ([xshift=0.45cm]current page.north west);
    \node[anchor=west, text=white, font=\huge\bfseries, text width=13cm]
      at ([xshift=1.1cm, yshift=1.2cm]current page.west) {Pixels, Frames \& Transformers};
    \node[anchor=west, text=wurgreen, font=\large, text width=13cm]
      at ([xshift=1.1cm, yshift=0.05cm]current page.west) {#1 · \coursemodule};
    \node[anchor=west, text=midgray, font=\small, text width=12cm]
      at ([xshift=1.1cm, yshift=-1.5cm]current page.west)
      {Denise J. Roth, PhD Candidate\\
       Strategic Communication Group · Wageningen University \& Research};
  \end{tikzpicture}
\end{frame}
}}

% ── Small helpers ───────────────────────────────────────────────────────────
\newcommand{\keyterm}[1]{\textbf{\textcolor{darkgreen}{#1}}}
\newcommand{\muted}[1]{\textcolor{midgray}{#1}}
\newcommand{\credit}[1]{{\tiny\textcolor{midgray}{#1}}}
\newcommand{\notebookslide}[1]{%
\begin{frame}{In the Notebook}
\vspace{0.1cm}
#1
\end{frame}}
\newcommand{\nbitem}[2]{\item[\textcolor{wurgreen}{#1}] #2}


\begin{document}

\moduletitle{Images and Words: CLIP}

% ─────────────────────────────────────────────────────────────────────────────
\begin{frame}{Learning from Captions Instead of Labels}
\begin{columns}[T]
\begin{column}{0.5\textwidth}
\textbf{CLIP} \muted{(Radford et al., 2021)} learned from about 400 million images with the text
that came with them on the web.

\vspace{0.15cm}
\begin{itemize}
\item An image encoder (a ViT) and a text encoder
\item In each batch, every image must pick out its own caption among all captions, and every
      caption its own image
\item \keyterm{Contrastive} learning: pull matching pairs together, push the rest apart
\end{itemize}
\end{column}
\begin{column}{0.46\textwidth}
\centering
\begin{tikzpicture}[x=0.62cm, y=0.62cm]
\foreach \i in {0,...,4} {
  \foreach \j in {0,...,4} {
    \ifnum\i=\j
      \fill[wurgreen] (\j, -\i) rectangle ++(0.9, -0.9);
    \else
      \fill[lightgray] (\j, -\i) rectangle ++(0.9, -0.9);
    \fi
  }
  \node[left, font=\scriptsize] at (-0.1, -\i - 0.45) {image \the\numexpr\i+1\relax};
  \node[above, font=\scriptsize, rotate=60, anchor=south west] at (\i + 0.3, 0.1) {caption \the\numexpr\i+1\relax};
}
\end{tikzpicture}\\[2pt]
\muted{\footnotesize the diagonal should score high, everything else low}
\end{column}
\end{columns}
\end{frame}

% ─────────────────────────────────────────────────────────────────────────────
\begin{frame}{One Space for Images and Sentences}
\begin{columns}[T]
\begin{column}{0.5\textwidth}
Both encoders map into the same space (512 numbers for CLIP ViT-B/32). Normalize every vector to
length 1, and the dot product of an image and a sentence is their \keyterm{cosine similarity}.

\vspace{0.2cm}
This lets you:
\begin{itemize}
\item classify with words (\keyterm{zero-shot})
\item search images with a sentence
\item define a measurement scale with two sentences
\end{itemize}
\end{column}
\begin{column}{0.46\textwidth}
\begin{warnbox}
\textbf{Similarities are relative.} Values usually fall between about 0.15 and 0.35, and a few
hundredths matter. Compare one image with several sentences, or one sentence with several images.
A similarity of 0.27 does not mean ``27\% protest''.
\end{warnbox}
\end{column}
\end{columns}
\end{frame}

% ─────────────────────────────────────────────────────────────────────────────
\begin{frame}{Zero-Shot Classification}
\vspace{0.1cm}
\centering
\begin{tikzpicture}[
  b/.style={draw=midgray, fill=lightgray, rounded corners=3pt, font=\footnotesize, align=center, inner sep=5pt},
  g/.style={draw=wurgreen, fill=lightgreen, rounded corners=3pt, font=\footnotesize, align=center, inner sep=5pt},
  arr/.style={-{Stealth[length=2mm]}, thick, color=midgray}]
\node[b, text width=2.4cm] (img) at (0, 0) {photo\\ $\rightarrow$ image encoder};
\node[b, text width=5.2cm] (t1) at (5.6, 0.9) {``a photo of a street festival''\\ $\rightarrow$ text encoder};
\node[b, text width=5.2cm] (t2) at (5.6, -0.9) {``a photo of a political protest''\\ $\rightarrow$ text encoder};
\node[g, text width=2.6cm] (out) at (11.0, 0) {most similar sentence wins\\ \muted{softmax of 100 $\times$ similarity}};
\draw[arr] (img.east) -- (t1.west);
\draw[arr] (img.east) -- (t2.west);
\draw[arr] (t1.east) -- (out.north west);
\draw[arr] (t2.east) -- (out.south west);
\end{tikzpicture}

\vspace{0.4cm}
\raggedright
No labeled examples needed: you describe each category in words. Every choice of words is a
measurement decision, so it needs the same scrutiny as a codebook definition.
\end{frame}

% ─────────────────────────────────────────────────────────────────────────────
\begin{frame}{The Wording Changes the Result}
\begin{columns}[T]
\begin{column}{0.5\textwidth}
\begin{itemize}
\item A few changed words can move accuracy by ten points or more, and the share labeled
      ``protest'' even more
\item Negation fails: ``a photo that does not show a protest'' mostly registers \emph{protest}
\item \keyterm{Prompt ensembles} average several templates (``a photo of \ldots'', ``a news photo
      of \ldots'') so no single wording dominates
\end{itemize}
\end{column}
\begin{column}{0.46\textwidth}
\begin{greenbox}
\textbf{Choose the wording on data you don't score on}\\
Picking the best prompt on the test set is tuning on the test set. Compare wordings on training
images (or a separate validation sample), then score once.
\end{greenbox}
\begin{graybox}
\textbf{Report the spread} across reasonable wordings, not just the winner.
\end{graybox}
\end{column}
\end{columns}
\end{frame}

% ─────────────────────────────────────────────────────────────────────────────
\begin{frame}{Zero-Shot or Labels?}
\begin{columns}[T]
\begin{column}{0.48\textwidth}
\textbf{A linear probe on the same embeddings} uses your labeled images. With a few hundred
labels, it usually beats zero-shot by a clear margin.

\vspace{0.2cm}
\textbf{Zero-shot is useful when}
\begin{itemize}
\item you have no labels yet
\item you need many categories quickly
\item you want candidates to hand-code
\end{itemize}
\end{column}
\begin{column}{0.48\textwidth}
\begin{graybox}
\textbf{SigLIP 2} \muted{(Google, 2025, Apache 2.0)}\\
A newer image-text model. It scores each image-caption pair on its own instead of against the
whole batch, and it is usually stronger than the original CLIP. Same workflow, different
\texttt{model\_name}.
\end{graybox}
\muted{\footnotesize A prompt chosen for one model is not necessarily best for another.}
\end{column}
\end{columns}
\end{frame}

% ─────────────────────────────────────────────────────────────────────────────
\begin{frame}{Searching a Collection with Words}
\begin{columns}[T]
\begin{column}{0.5\textwidth}
Embed every image once; then any sentence is a query.

\vspace{0.15cm}
\begin{itemize}
\item ``packaging with a cartoon character''
\item ``a chocolate cereal for children''
\end{itemize}

\vspace{0.15cm}
\keyterm{Image + words}: add the embedding of a phrase to the embedding of an image (``this box,
\emph{with a cartoon character}'') to find similar items that differ in one respect.
\end{column}
\begin{column}{0.46\textwidth}
\begin{graybox}
\textbf{Good for}\\
getting into a large collection, finding examples and counterexamples for a codebook, and
building hand-coding samples that are rich in rare cases.
\end{graybox}
\begin{warnbox}
Search results are suggestions to look at, not measurements.
\end{warnbox}
\end{column}
\end{columns}
\end{frame}

% ─────────────────────────────────────────────────────────────────────────────
\begin{frame}{A Measurement Scale Defined in Words}
\begin{columns}[T]
\begin{column}{0.5\textwidth}
Score each image by its similarity to one pole minus its similarity to the other:

\vspace{0.1cm}
{\footnotesize
\textcolor{darkgreen}{``a wholesome, natural breakfast cereal with whole grains''}\\
\hspace{1.5cm}versus\\
\textcolor{warnborder}{``a sweet, sugary breakfast cereal for children, a fun treat''}}

\vspace{0.2cm}
Like the poles in the embedding playground of Bags, Vectors \& Transformers, but for images.
\end{column}
\begin{column}{0.46\textwidth}
\begin{greenbox}
\textbf{Validate it}\\
Against human ratings of the concept you mean. Nutri-Score measures something related but
different: what is \emph{in} the box, not what the box \emph{suggests}. Mismatches can be the
finding.
\end{greenbox}
\begin{graybox}
\textbf{Check the wording}: rebuild the scale with other poles and see how much items move.
\end{graybox}
\end{column}
\end{columns}
\end{frame}

% ─────────────────────────────────────────────────────────────────────────────
\begin{frame}{What CLIP Learned About People}
\begin{columns}[T]
\begin{column}{0.5\textwidth}
CLIP learned from web captions, with their habits about who does what. Its own authors found that
some groups in FairFace were more often put into crime-related or non-human categories
\muted{(Radford et al., 2021, section 7)}.

\vspace{0.2cm}
\textbf{An audit with a known answer}: every Congress portrait shows a politician. How often does
CLIP pick ``a politician'' from a list of occupations, for women and for men?
\end{column}
\begin{column}{0.46\textwidth}
\begin{warnbox}
\textbf{The label set shapes the result}\\
Add ``a member of Congress'' or remove ``a nurse'' and ``a secretary'', and both the share correct
and the gap between women and men change.
\end{warnbox}
\muted{\footnotesize Report the full label set, and check that conclusions hold across reasonable
alternatives (compare Schwemmer et al., 2020, in module 07).}
\end{column}
\end{columns}
\end{frame}

% ─────────────────────────────────────────────────────────────────────────────
\notebookslide{
\begin{itemize}
\setlength{\itemsep}{5pt}
\nbitem{\faComments}{Images and sentences in one space}
\nbitem{\faTags}{Zero-shot classification of your 40 coded images, with prompts chosen on training data}
\nbitem{\faBalanceScale}{Zero-shot versus a linear probe, and SigLIP 2 as an alternative}
\nbitem{\faSearch}{Searching the packaging with words, and with an image plus words}
\nbitem{\faRuler}{A ``healthy-looking'' scale, checked against Nutri-Score}
\nbitem{\faUsers}{An occupation audit on the Congress portraits}
\end{itemize}

\vspace{0.35cm}
\muted{\footnotesize part3\_transformers/09\_clip\_images\_and\_words/notebooks/\\
09\_clip\_images\_and\_words\_exercises.ipynb · 6 exercises · GPU recommended}
}

\end{document}
